% Source: https://artofproblemsolving.com/community/c3h1521666_need_help
% Stable identity: aops-topic-1521666-q01
The ten lines give $20$ rays from the point, and consecutive rays (going around the point) form $20$ angles that together make up the full angle (one of them is marked in red):

\begin{asy}
usepackage("amsmath");
size(4.5cm);
defaultpen(fontsize(10pt));
real[] ang={0,14,40,57,75,95,118,133,150,168};
for(int i=0;i<10;++i){ draw(2.2*dir(ang[i])--2.2*dir(ang[i]+180),linewidth(0.8)); }
dot((0,0));
draw(arc((0,0),0.9,0,14),red+linewidth(0.9));
label("\mbox{}",(min(currentpicture,user=true)+max(currentpicture,user=true))/2,black); // keeps the text after the figure black when LaTeX typesets the labels inline
\end{asy}

These $20$ angles add up to $360^\circ$. If every one of them were at least $19^\circ$, their sum would be at least $20\cdot19^\circ=380^\circ>360^\circ$, which is impossible. Hence some angle between neighbouring rays, which is an angle between two of the lines, is at most $\dfrac{360^\circ}{20}=18^\circ<19^\circ$. $\blacksquare$

(In fact $18^\circ$ is best possible: ten lines through the point at $18^\circ$ apart make all twenty angles exactly $18^\circ$.)
